%!TEX root = ../template.tex
% Comentarios da revisao critica de 2026-09-13, visiveis no PDF.
% Para a versao sem comentarios, trocar a linha seguinte por
% \mostrarrevisaofalse e recompilar template.tex.
\newif\ifmostrarrevisao
\mostrarrevisaotrue

\tcbuselibrary{breakable}
\newcommand{\revisaotese}[5]{%
  \ifmostrarrevisao
    \par\begingroup
    \colorlet{revisaocor}{blue!55!black}%
    \ifstrequal{#2}{P1}{\colorlet{revisaocor}{red!65!black}}{}%
    \ifstrequal{#2}{P3}{\colorlet{revisaocor}{black!60}}{}%
    \ifstrequal{#3}{FIGURA}{\colorlet{revisaocor}{teal!65!black}}{}%
    \begin{tcolorbox}[
      breakable,
      colback=revisaocor!4!white,
      colframe=revisaocor,
      colbacktitle=revisaocor,
      coltitle=white,
      boxrule=0.5pt,arc=1mm,
      left=2.5mm,right=2.5mm,top=2mm,bottom=2mm,
      before skip=8pt,after skip=8pt,
      fonttitle=\sffamily\bfseries\small,
      fontupper=\normalfont\small,
      title={REVISÃO #1 | #2 | #3}]
      \setlength{\emergencystretch}{3em}%
      \textbf{#4}\par\smallskip
      #5
    \end{tcolorbox}%
    \endgroup
  \fi
}

% As definicoes de glossarios sao lidas no preambulo. As tres notas
% sobre esses ficheiros aparecem neste apendice apenas em modo de revisao.
\ifmostrarrevisao
  % 2026-09-18: Apendice H arquivado em _revisao/ (so tinha FR-01/02/03,
  % resolvidas a 2026-09-17 com as siglas, simbolos e glossario reais).
  % \ntaddfile{appendix}{appendix-H-revisao}
\fi

% ---------------------------------------------------------------------
% 2026-09-19: marcacao temporaria das alteracoes da revisao integral.
% \alterado{...} destaca o texto novo ou reescrito. Para retirar TODAS as
% marcas de uma vez, trocar a definicao activa pela neutra, logo abaixo.
% (Nao se usa soul/\hl: rebenta com acentos, \ref e \cite em pdflatex.)
\definecolor{corAlterado}{RGB}{176,0,132}
%\DeclareRobustCommand{\alterado}[1]{\begingroup\color{corAlterado}#1\endgroup}
\DeclareRobustCommand{\alterado}[1]{#1}   %% 2026-09-21: camada fechada, sem marcas
% ---------------------------------------------------------------------

% ---------------------------------------------------------------------
% 2026-09-19: marcacao do material NOVO acrescentado depois da passagem
% de correccoes (Secao 2.3.3, comparadores industriais e legais). Ciano,
% para se distinguir do magenta das correccoes. Mesma logica de remocao.
\definecolor{corNovo}{RGB}{0,122,163}
%\DeclareRobustCommand{\novo}[1]{\begingroup\color{corNovo}#1\endgroup}
\DeclareRobustCommand{\novo}[1]{#1}   %% 2026-09-21: camada fechada, sem marcas
% ---------------------------------------------------------------------

% ---------------------------------------------------------------------
% 2026-09-20: marcacao da passagem de revisao integral do corpo (Caps. 3-8),
% a partir da analise externa em _revisao/2026-09-20-corpo-tese-analise.md.
% Roxo, para se distinguir do magenta (correccoes do Cap. 2) e do ciano
% (material novo). Mesma logica de remocao.
\definecolor{corRevisto}{RGB}{126,47,186}
%\DeclareRobustCommand{\revisto}[1]{\begingroup\color{corRevisto}#1\endgroup}
\DeclareRobustCommand{\revisto}[1]{#1}   %% 2026-09-21: camada fechada, sem marcas
% ---------------------------------------------------------------------

%% ---------------------------------------------------------------------
%% Camada de fecho (2026-09-21). Laranja queimado: e a unica marca quente
%% do documento, distinta das tres camadas anteriores (magenta, ciano,
%% roxo) e da cor das referencias cruzadas. Para a versao final, trocar
%% as duas linhas abaixo.
\definecolor{corFecho}{RGB}{181,69,11}
%\DeclareRobustCommand{\fecho}[1]{\begingroup\color{corFecho}#1\endgroup}
\DeclareRobustCommand{\fecho}[1]{#1}   %% 2026-09-21: camada fechada, sem marcas

%% ---------------------------------------------------------------------
%% Camada de correccao (2026-09-21, tarde). Magenta, a mesma cor da
%% primeira passagem de correccoes do Cap. 2 -- reaberta aqui porque o
%% que marca e da mesma natureza: numeros, denominadores e afirmacoes
%% corrigidos depois da analise critica do Cap. 3. Para a versao final,
%% trocar as duas linhas abaixo.
\DeclareRobustCommand{\corr}[1]{\begingroup\color{corAlterado}#1\endgroup}
%\DeclareRobustCommand{\corr}[1]{#1}   %% <- versao final, sem marcas

%% ---------------------------------------------------------------------
%% Camada de ajuste (2026-09-21, fim de tarde). Ciano, a cor do material
%% novo -- reaberta para a passagem sobre a Seccao 3.5: jargao de
%% instrumento traduzido, duplicacoes cortadas e leitura de engenharia
%% acrescentada. Distinta do magenta, que marca as correccoes de numeros
%% nas Seccoes 3.1 a 3.3. Para a versao final, trocar as duas linhas.
\DeclareRobustCommand{\ajuste}[1]{\begingroup\color{corNovo}#1\endgroup}
%\DeclareRobustCommand{\ajuste}[1]{#1}   %% <- versao final, sem marcas
